\documentclass[11pt,a4paper]{article}
\usepackage[margin=2.4cm]{geometry}
\usepackage{amsmath,amssymb,amsthm,mathtools}
\usepackage{booktabs,array,float}
\usepackage[T1]{fontenc}
\usepackage{lmodern}
\usepackage{microtype}
\usepackage{xcolor}
\usepackage[colorlinks=true,linkcolor=blue!50!black,urlcolor=blue!50!black,citecolor=blue!50!black]{hyperref}

\newtheorem{theorem}{Theorem}[section]
\newtheorem{lemma}[theorem]{Lemma}
\newtheorem{proposition}[theorem]{Proposition}
\newtheorem{corollary}[theorem]{Corollary}
\theoremstyle{definition}
\newtheorem{definition}[theorem]{Definition}
\newtheorem{example}[theorem]{Example}
\theoremstyle{remark}
\newtheorem{remark}[theorem]{Remark}

\newcommand{\F}{\mathbb{F}}
\newcommand{\Z}{\mathbb{Z}}
\newcommand{\Q}{\mathbb{Q}}
\newcommand{\Fpbar}{\overline{\F}_p}
\newcommand{\leg}[2]{\left(\frac{#1}{#2}\right)}
\newcommand{\Hl}{\mathit{Hl}}
\newcommand{\Res}{\operatorname{Res}}
\newcommand{\disc}{\operatorname{disc}}
\newcommand{\ssp}{\mathit{ss}_p}
\newcommand{\Zp}{\Z_{(p)}}
\newcommand{\cL}{\mathcal{L}}
\newcommand{\cA}{\mathcal{A}}
\newcommand{\cB}{\mathcal{B}}
\newcommand{\cC}{\mathcal{C}}
\newcommand{\dw}{\partial_w}

\numberwithin{equation}{section}
\hypersetup{pdftitle={Fricke supersingular polynomials at levels 5, 7 and 13},pdfauthor={David Jordan}}
\setlength{\emergencystretch}{1.5em}

\title{Fricke supersingular polynomials at levels 5, 7 and 13}
\author{David Jordan}
\date{14 November 2026}

\begin{document}
\maketitle

\begin{abstract}
For $N\in\{5,7,13\}$ and every prime $p\ge5$ with $p\ne N$, we describe the
supersingular polynomial of the Fricke group $\Gamma_0^*(N)$ as a product of
explicit elliptic and Fricke factors and a polynomial solution, modulo $p$, of a
second-order differential equation; the coefficients of that solution are given
by a recurrence. At levels $5$ and $7$ the equation is of Heun type. The
result proves a conjecture of Sakai and its extension to all primes by Nakaya
(his Conjecture~1), once the first Heun parameter in Nakaya's formulation is
corrected; as printed there, the formula holds only for $p\equiv1\pmod4$ at
level $5$ and for $p\equiv1\pmod3$ and $p=5$ at level $7$. The same
argument proves Nakaya's Conjecture~3 on the squares of these polynomials. At
level $13$ the equation has five singular points and the recurrence has five
terms. One proof covers the three levels: the Hasse invariant, lifted by
$E_{p-1}$, is descended through the Fricke involution and compared with the
uniformizing differential equation of $X_0(N)$.
\end{abstract}

\section{Introduction}\label{sec:intro}

For a prime $p$, the supersingular polynomial $\ssp(J)\in\F_p[J]$ is the monic
polynomial whose roots are the $j$-invariants of the supersingular elliptic
curves in characteristic $p$. Classically it is a truncated hypergeometric
series, and Kaneko and Zagier~\cite{KZ} put this and several other descriptions,
through modular forms of weight $p-1$ and through Atkin's orthogonal
polynomials, into one framework. Analogous polynomials for congruence subgroups
of low level were studied by Tsutsumi~\cite{Tsutsumi}. For the Fricke groups
$\Gamma_0^*(N)=\langle\Gamma_0(N),w_N\rangle$, a supersingular polynomial
$\ssp^{(N*)}$ was introduced for $N\le7$ by Koike and
Sakai (see~\cite{Sakai11,Sakai14}) and, for each prime $N$ for which
$X_0(N)/w_N$ has genus zero, by Nakaya~\cite{Nakaya19}.

For $N=5$ and $7$, Sakai proved that the Eisenstein series $f_N$ of weight two
is the square of a Heun series in $1/j_N^*$, and conjectured that, for primes
$p$ with $\leg{-N}{p}=1$, the same Heun series gives $\ssp^{(N*)}$ modulo
$p$~\cite[Proposition~3.1 and Conjecture~5.2]{Sakai14}. Nakaya's
Conjecture~1~\cite[pp.~488--489]{Nakaya19} restates this with the first Heun
parameter written as a negative integer and extends it to all primes $p\ne N$,
and his Conjecture~3 expresses the square of $\ssp^{(N*)}$ through Ap\'ery-like
numbers.

\subsection*{Results}
This paper treats the three levels $N=5,7,13$. Together with the
hypergeometric levels $2$ and $3$, these are the prime levels for which
$X_0(N)$ itself has genus zero, and the proof uses a Hauptmodul of $X_0(N)$.

\begin{itemize}
\item Theorem~\ref{thm:uniform} gives, for each of the three levels and every
prime $p\ge5$, $p\ne N$, a factorization of the supersingular polynomial into
factors coming from the elliptic points, a factor coming from the Fricke fixed
points, and a polynomial $Q_p$. The reciprocal polynomial of $Q_p$ is the
unique normalized polynomial solution of degree $<p$ of an explicit
second-order differential equation modulo $p$, and its coefficients are
generated by a recurrence.
\item At levels $5$ and $7$ the recurrence has three terms
(Theorem~\ref{thm:rat}) and the equation is Heun's. Theorem~\ref{thm:main}
proves Sakai's conjecture as he stated it (Corollary~\ref{cor:sakai}) and
extends it to all primes. In Nakaya's formulation this requires one change:
the first Heun parameter must be
\[
\varepsilon_p-m_5\qquad\text{and}\qquad\delta_p-m_7
\]
instead of the printed $-m_5$ and $-m_7$, where
$\varepsilon_p=\frac12(1-\leg{-1}{p})$ and $\delta_p=\frac12(1-\leg{-3}{p})$.
The printed formula is correct exactly when the elliptic point of the Fricke
quotient is not supersingular, and at one further prime
(Corollary~\ref{cor:printed}). The smallest failure is $N=7$, $p=11$, a prime
at which Sakai's original statement applies and is correct.
Corollary~\ref{cor:conj3} proves Nakaya's Conjecture~3.
\item At level $13$ the quotient has two elliptic points, the equation has
five singular points, and the recurrence has five terms
(Theorem~\ref{thm:13}). There are eight cases, indexed by the characters
$\leg{-1}{p}$, $\leg{-3}{p}$ and $\leg{-13}{p}$. A new feature is that the
recurrence polynomial can have one degree less than expected; this happens
exactly when the curve above the base point of the coordinate is supersingular,
first at $p=149$ (Theorem~\ref{thm:13}(iii) and Remark~\ref{rem:basepoint}).
\end{itemize}

\subsection*{Method}
The method is the modular-form method of~\cite{KZ}, and it is the same at the
three levels. Let $h$ be the eta-quotient Hauptmodul of $X_0(N)$ and
$f_N=(NE_2(N\tau)-E_2(\tau))/(N-1)$ the Eisenstein series of weight two.
\begin{enumerate}
\item By Deligne's congruence~\cite{Serre}, \cite[Theorem~1]{KZ}, the form $E_{p-1}$ lifts
the Hasse invariant. Dividing by $f_N^{(p-1)/2}$ gives a rational function of
$h$ whose numerator vanishes modulo $p$ at the supersingular points other
than the elliptic points
(Proposition~\ref{prop:Np}).
\item Since $E_{p-1}\equiv1\pmod p$, this numerator is, modulo $p$, invariant
or anti-invariant under the Fricke involution $h\mapsto K/h$, with sign
$\leg{-N}{p}$. It therefore descends to a polynomial $Q_p$ in the Fricke
coordinate $x=h+A+K/h$ (Lemma~\ref{lem:laurent}).
\item The series $G=f_N^{1/2}$ in $w=1/x$ satisfies a second-order
differential equation (Lemma~\ref{lem:unif}). The same congruence gives
$G(w)\equiv P_p(w)\,G(w^p)$, up to elementary factors, where $P_p$ is the
reciprocal of $Q_p$. Because differentiation kills $G(w^p)$ modulo $p$, the
polynomial $P_p$ satisfies the equation modulo $p$. That a Hasse invariant is
a truncated solution of a Picard--Fuchs equation in this way goes back to
Igusa~\cite{Igusa} and Dwork~\cite{Dwork}.
\item A polynomial solution of degree $<p$ is determined by its constant
term. This identifies $P_p$ with the output of the recurrence and proves, in
particular, that the recurrence terminates.
\end{enumerate}
Termination is a consequence of the argument, not an input: a non-positive
integral Heun parameter does not by itself make a Heun series terminate, since
the accessory parameter must also be special. The elliptic and Fricke factors
have a uniform meaning. An elliptic point of order $2$ or $3$ contributes a
linear factor when $j=1728$ or $j=0$ is supersingular, and the Fricke fixed
points contribute their quadratic polynomial when the curves with complex
multiplication by $\Z[\sqrt{-N}]$ are supersingular. At levels $5$ and $7$ the
single elliptic point is the base point $x=0$ of the coordinate, which is why
its contribution appears as a power of $X$ and why the equation has only four
singular points.

\subsection*{Related work}
\emph{Sakai and Nakaya.} The Heun identity for $f_5$ and $f_7$ and the
characteristic-zero recurrences for its coefficients are due to
Sakai~\cite[Proposition~3.1 and Lemma~3.3]{Sakai14}. They are re-derived here
by a different argument (Proposition~\ref{prop:heun}); the operators for
$\mu_N=1$ and the congruences modulo $p$ are not in~\cite{Sakai14}. Sakai's
Conjecture~5.2 concerns the primes with $\leg{-N}{p}=1$ and uses the Heun
parameters $(\frac14,\frac34)$ and $(\frac13,\frac23)$; it is correct as
stated. The discrepancy described in Corollary~\ref{cor:printed} first appears
in the reformulation of~\cite[Conjecture~1]{Nakaya19}; see
Remark~\ref{rem:typo}.

\emph{Other conjectures of Nakaya.} At levels $5$ and $7$,
Morton~\cite{Morton23} proved Conjecture~4 of~\cite{Nakaya19}, a resultant
formula for $(\ssp^{(N*)})^2$, together with the degree of $\ssp^{(N*)}$, and
Conjecture~2, which says that $\ssp^{(N*)}$ splits into linear factors exactly
for the primes dividing the order of the Harada--Norton group ($N=5$) or the
Held group ($N=7$). The latter was also obtained by Aricheta~\cite{Aricheta}
in the setting of monstrous moonshine. Morton~\cite{Morton21a} showed that the
roots of $\ssp^{(N*)}$ lie in $\F_{p^2}$, and the number of linear factors
(Conjecture~5) is the subject of~\cite{Morton21a,Morton21b,Morton24}; the
introduction of~\cite{Morton24} records the Heun description as a conjecture
of Sakai and Nakaya. Conjectures~5 and~6 of~\cite{Nakaya19} are treated for all
levels in~\cite{JordanNakaya}.

\emph{Ap\'ery-like numbers.} Lucas congruences for the sequences $u_5$, $u_7$
of Conjecture~3 are known~\cite{MalikStraub,BTY}; their relation to
Corollary~\ref{cor:conj3} is explained in Remark~\ref{rem:lucas}.

\emph{Level $13$.} The parametrization of $j$ by $h$ goes back to
Klein~\cite{Klein}; see Maier~\cite{Maier} for a modern account of the
genus-zero curves $X_0(N)$ and their differential equations. A third-order
differential equation for an Eisenstein series of weight two and level $13$
was given by Cooper and Ye~\cite{CooperYe}.

\emph{General theory.} Relations between Hasse invariants, Picard--Fuchs
equations and Atkin polynomials for genus-zero curves in Hilbert modular
varieties are developed in~\cite{Bogo,BL1,BL2}. These results allow elliptic
points: \cite[Corollary~3]{Bogo} describes non-ordinary loci by truncating a
power of the Picard--Fuchs solution, and \cite{BL1} constructs lifts of
partial Hasse invariants under a condition on the elliptic
points~\cite[(20)]{BL1} which, as we read it, fails for the three Fricke
groups considered here (each has several elliptic points of order two), with
a variant that does not need it~\cite[Remark~5]{BL1}. We have
not tried to deduce the present formulas from that theory, and we do not know
to what extent it implies them.

To the author's knowledge, a proof of Sakai's Conjecture~5.2, a proof or
correction of Conjectures~1 and~3 of~\cite{Nakaya19}, and the level-$13$
recurrence have not appeared in print. The search behind this statement cannot
be exhaustive, and no claim of priority is made for a method that is classical.

\subsection*{Organization}
Section~\ref{sec:results} fixes the notation and states the results.
Section~\ref{sec:proof} proves Theorem~\ref{thm:uniform} for the three levels
at once. Section~\ref{sec:57} derives the Heun form and the two corollaries at
levels $5$ and $7$, and Section~\ref{sec:13} discusses level $13$. The proofs
use finitely many identities between explicit modular forms and polynomials;
Appendix~\ref{app:certificates} describes the exact computations that certify
them.

\section{Statement of results}\label{sec:results}

\subsection{Setting}\label{subsec:setting}

Throughout, $N\in\{5,7,13\}$ and $p\ge5$ is a prime with $p\ne N$. Put
\[
\varepsilon_p=\frac{1-\leg{-1}{p}}2,\qquad
\delta_p=\frac{1-\leg{-3}{p}}2,\qquad
\mu_N=\frac{1-\leg{-N}{p}}2 ,
\]
so that $\varepsilon_p=1$, $\delta_p=1$, $\mu_N=1$ exactly when the curves
with $j=1728$, with $j=0$, and with complex multiplication by $\Z[\sqrt{-N}]$,
respectively, are supersingular in characteristic $p$. We often write
$\varepsilon,\delta,\mu$.

Let $q=e^{2\pi i\tau}$ and let $\eta$ be the Dedekind eta function. The
function
\[
h=h_N=\Bigl(\frac{\eta(\tau)}{\eta(N\tau)}\Bigr)^{24/(N-1)}=q^{-1}+\cdots
\]
is a Hauptmodul of the genus-zero curve $X_0(N)$, and the Fricke involution
$\tau\mapsto-1/N\tau$ sends $h$ to $K/h$ with $K=N^{12/(N-1)}$
(Lemma~\ref{lem:basic}). We use the Fricke-invariant coordinate
\[
x=h+A+\frac Kh ,\qquad w=\frac1x ,
\]
with the constants of Table~\ref{tab:data}. At levels $5$ and $7$ we write
$X$ for $x$; this is Nakaya's Hauptmodul $j_N^*$~\cite[p.~485]{Nakaya19}. At
level $13$ we write $t=h+13/h$; since $t=q^{-1}-2+O(q)$ and Nakaya normalizes
$j_{13}^*$ to have constant term $0$, we have $j_{13}^*=t+2$. The coordinate
$t$ makes the Fricke polynomial even.

\begin{table}[htbp]
\centering\small
\renewcommand{\arraystretch}{1.35}\setlength{\tabcolsep}{6pt}
\begin{tabular}{@{}llll@{}}
\toprule
 & $N=5$ & $N=7$ & $N=13$\\
\midrule
$h$ & $(\eta(\tau)/\eta(5\tau))^6$ & $(\eta(\tau)/\eta(7\tau))^4$ & $(\eta(\tau)/\eta(13\tau))^2$\\
$K,\ A$ & $125,\ 22$ & $49,\ 13$ & $13,\ 0$\\
coordinate $x=h+A+K/h$ & $X=j_5^*$ & $X=j_7^*$ & $t=j_{13}^*-2$\\
$d_2(h)$ & $h^2+22h+125$ & $1$ & $h^2+6h+13$\\
$d_3(h)$ & $1$ & $h^2+13h+49$ & $h^2+5h+13$\\
elliptic polynomial $\xi_N(x)$ & $X$ & $X$ & $(t+6)(t+5)$\\
elliptic factor $\Xi_{N,p}(x)$ & $X^{\varepsilon_p}$ & $X^{\delta_p}$ & $(t+6)^{\varepsilon_p}(t+5)^{\delta_p}$\\
Fricke polynomial $B_N(x)$ & $X^2-44X-16$ & $(X+1)(X-27)$ & $t^2-52$\\
$s(w)=w^2B_N(1/w)$ & $1-44w-16w^2$ & $1-26w-27w^2$ & $1-52w^2$\\
$\ell_2(w),\ \ell_3(w)$ & $1,\ 1$ & $1,\ 1$ & $1+6w,\ 1+5w$\\
$e$ & $e_2$ & $e_3$ & $e_2+e_3$\\
\bottomrule
\end{tabular}
\caption{Level data. The polynomials $d_2$, $d_3$ vanish at the elliptic points
of $X_0(N)$ of order $2$ (where $j=1728$) and of order $3$ (where $j=0$), and
$d_i=1$ when there is none; $d_i(h)/h$ is the corresponding linear factor of
$\xi_N(x)$. The exponents $e_2$, $e_3$ are defined in~\eqref{eq:exps}.}
\label{tab:data}
\end{table}

The polynomial $B_N(x)=(x-A)^2-4K$ vanishes at the images of the fixed points
$h^2=K$ of the Fricke involution. The elliptic polynomial $\xi_N$ vanishes at
the images of the elliptic points: $X=0$ at levels $5$ and $7$, and $t=-6$
(order $2$) and $t=-5$ (order $3$) at level $13$.

\paragraph{Supersingular polynomials.}
Let $\ssp(J)\in\F_p[J]$ be the monic polynomial whose roots are the
supersingular $j$-invariants in $\Fpbar$, each taken once. Let
$R_N(J,x)=J^2-a_N(x)J+b_N(x)\in\Z[J,x]$ be the modular relation with
$R_N(J,x(\tau))=(J-j(\tau))(J-j(N\tau))$. For $N=5,7$ these are the
polynomials of~\cite[Example~1]{Nakaya19},
\begin{align*}
R_5(J,X)&=J^2-J\bigl(X^5-80X^4+1890X^3-12600X^2+7776X+3456\bigr)
              +(X^2+216X+144)^3,\\
R_7(J,X)&=J^2-JX(X^2-21X+8)(X^4-42X^3+454X^2-1008X-1280)\\
&\hspace{22mm}+X^2(X^2+224X+448)^3,
\end{align*}
and for $N=13$, in the coordinate $t$,
\begin{align*}
a_{13}(t)&=(t+5)(t^4+7t^3-32t^2-254t-180)\\
&\quad\times(t^8+14t^7-15t^6-956t^5-2895t^4+12974t^3+60805t^2-10836t-194076),\\
b_{13}(t)&=(t+5)^2(t^4+254t^3+5077t^2+34092t+75492)^3
\end{align*}
(Lemma~\ref{lem:ids}(e)).

\begin{definition}\label{def:S}
$S_{N,p}(x)\in\F_p[x]$ is the monic squarefree polynomial with the same roots
in $\Fpbar$ as $\Res_J\bigl(\ssp(J),R_N(J,x)\bigr)$.
\end{definition}

Thus $x_0$ is a root of $S_{N,p}$ exactly when one of the two $N$-isogenous
curves above $x_0$ is supersingular, and then both are. The definition involves no
differential equation. For $N=5,7$ it is Nakaya's
$\ssp^{(N*)}(X)$~\cite[eq.~(20)]{Nakaya19}, and
$\ssp^{(13*)}(X)=S_{13,p}(X-2)$.

\paragraph{Exponents.} Put
\begin{equation}\label{eq:exps}
e_2=\frac{p-1-2\varepsilon_p}{4}=\Bigl\lfloor\frac{p-1}4\Bigr\rfloor,\qquad
e_3=\frac{p-1-\delta_p}{3}=\Bigl\lfloor\frac{p-1}3\Bigr\rfloor,\qquad
n=e-\mu_N,
\end{equation}
with $e$ as in Table~\ref{tab:data}, and
\begin{equation}\label{eq:rhos}
\rho_2=-\frac{1+2\varepsilon_p}{4},\qquad\rho_3=-\frac{1+\delta_p}{3},
\qquad\text{so that}\quad e_2-\rho_2=\frac p4,\quad e_3-\rho_3=\frac p3 .
\end{equation}

\subsection{The differential operators}\label{subsec:ops}

Let
\begin{equation}\label{eq:Psi}
\Psi_5(w)=3+3w,\qquad\Psi_7(w)=2+6w,\qquad
\Psi_{13}(w)=\frac{1+23w+165w^2+381w^3}{(1+6w)^2(1+5w)^2}.
\end{equation}
The \emph{uniformizing operator} of level $N$ is
\begin{equation}\label{eq:LG}
L_N=ws\,\dw^2+\bigl(s+\tfrac12ws'\bigr)\dw-\Psi_N(w),\qquad s=s(w),\quad \dw=\frac{d}{dw}.
\end{equation}
It annihilates the series $G=f_N^{1/2}$ in $w$ (Lemma~\ref{lem:unif}). The
operator that governs the supersingular polynomial is its conjugate by an
elementary factor. Put
\begin{equation}\label{eq:Theta}
\Theta=\frac{d}{dw}\log\bigl(\ell_2^{\rho_2}\,\ell_3^{\rho_3}\,s^{-\mu/2}\bigr)
=\rho_2\frac{\ell_2'}{\ell_2}+\rho_3\frac{\ell_3'}{\ell_3}-\frac{\mu s'}{2s}
\end{equation}
and define $\cL_{N,p}=\cA\,\dw^2+\cB\,\dw+\cC$ by
\begin{equation}\label{eq:gaugeop}
\begin{aligned}
\cA&=ws\,\ell_2\ell_3,\\
\cB&=\ell_2\ell_3\bigl(s+\tfrac12ws'-2ws\,\Theta\bigr),\\
\cC&=\ell_2\ell_3\bigl(ws(\Theta^2-\Theta')-(s+\tfrac12ws')\Theta-\Psi_N\bigr).
\end{aligned}
\end{equation}
The operator depends on $p$ only through $(\varepsilon_p,\delta_p,\mu_N)$, and
at levels $5$ and $7$, where $\ell_2=\ell_3=1$, only through $\mu_N$.
Lemma~\ref{lem:gauge} shows that $\cA,\cB,\cC$ are polynomials in $w$ with
coefficients in $\Z[\frac16]$, that $\cA(0)=0$ and $\cA'(0)=\cB(0)=1$, and that
\begin{equation}\label{eq:Lmu}
\cL_{N,p}=L_\mu:=ws\,\dw^2+\bigl(s+(\mu+\tfrac12)ws'\bigr)\dw-\Psi_{N,\mu}(w)
\qquad(N=5,7),
\end{equation}
with the linear polynomials $\Psi_{N,\mu}$ of Table~\ref{tab:rec}. The eight
operators of level $13$ are listed in Table~\ref{tab:ops13}.

Write $\cA=\sum_i\cA_iw^i$, and similarly for $\cB$, $\cC$. For a series
$\sum_jg_jw^j$, the coefficient of $w^j$ in $\cL_{N,p}\bigl(\sum g_iw^i\bigr)$ is
\begin{equation}\label{eq:rec}
(j+1)^2g_{j+1}+\sum_{i\ge2}\cA_i\,(j-i+2)(j-i+1)\,g_{j-i+2}
+\sum_{i\ge1}\cB_i\,(j-i+1)\,g_{j-i+1}+\sum_{i\ge0}\cC_i\,g_{j-i},
\end{equation}
because $\cA_0=0$ and $\cA_1=\cB_0=1$.

\subsection{The main theorem}

\begin{theorem}\label{thm:uniform}
Let $N\in\{5,7,13\}$ and let $p\ge5$ be a prime, $p\ne N$. Define $g_j\in\F_p$
for $j\le p-1$ by $g_j=0$ for $j<0$, $g_0=1$, and by requiring
that~\eqref{eq:rec} vanish in $\F_p$ for $0\le j\le p-2$.
\begin{enumerate}
\item[(i)] $g_j=0$ for $n<j\le p-1$. The polynomial $P_p(w)=\sum_{j=0}^{n}g_jw^j$
satisfies $\cL_{N,p}P_p=0$ in $\F_p[w]$, and it is the only polynomial of
degree $<p$ with constant term $1$ that does.
\item[(ii)] $Q_p(x)=x^nP_p(1/x)=\sum_{j=0}^{n}g_jx^{n-j}$ is monic of degree
$n$, squarefree, and prime to $\xi_N(x)B_N(x)$.
\item[(iii)] The supersingular polynomial is
\begin{equation}\label{eq:mainformula}
S_{N,p}(x)=\Xi_{N,p}(x)\,B_N(x)^{\mu_N}\,Q_p(x)
=\Xi_{N,p}(x)\,B_N(x)^{\mu_N}\,x^{n}P_p(1/x) .
\end{equation}
In particular $\deg S_{N,p}=n+2\mu_N+\deg\Xi_{N,p}$.
\item[(iv)] $g_n=Q_p(0)$. Hence $\deg P_p=n$ unless $x=0$ is a root of $Q_p$,
in which case $\deg P_p=n-1$.
\end{enumerate}
\end{theorem}

The three factors in~\eqref{eq:mainformula} are pairwise coprime, and each has
a meaning: $\Xi_{N,p}$ collects the supersingular elliptic points, $B_N^{\mu_N}$
the supersingular Fricke fixed points, and $Q_p$ all other supersingular
points of $X_0(N)/w_N$.

By (iii) and~\eqref{eq:exps}, the degree of $S_{N,p}$ is
$\frac{p-1}4+\frac{\varepsilon_p}2+\mu_5$ for $N=5$,
$\frac{p-1}3+\frac{2\delta_p}3+\mu_7$ for $N=7$, and
$\frac{7(p-1)}{12}+\frac{\varepsilon_p}2+\frac{2\delta_p}3+\mu_{13}$ for $N=13$.
These are the values predicted by~\cite[Conjecture~6]{Nakaya19}. For $N=5,7$
the degree is stated in~\cite[Remark~5.3]{Sakai14} and was proved by
Morton~\cite{Morton23}; for all levels of Nakaya's list, including $N=13$,
see~\cite{JordanNakaya}. The degree is therefore not new here, but the proof
is independent.

\subsection{Levels 5 and 7}

At these levels the elliptic point is $X=0$, the elliptic factor is $X^{o_p}$
with
\[
o_p=\varepsilon_p\ (N=5),\qquad o_p=\delta_p\ (N=7),
\]
and we use Nakaya's notation
\begin{equation}\label{eq:mN}
m_N=\frac{p-\chi_p}{r}-\mu_N,\qquad
(r,\chi_p)=\bigl(4,\leg{-1}{p}\bigr)\ (N=5),\quad
(r,\chi_p)=\bigl(3,\leg{-3}{p}\bigr)\ (N=7).
\end{equation}
Then $e=\lfloor(p-1)/r\rfloor=(p-\chi_p)/r-o_p$ and $n=m_N-o_p$. Write
$s(w)=1-Tw-Uw^2$, so $(T,U)=(44,16)$ and $(26,27)$.

\begin{table}[htbp]
\centering\small
\renewcommand{\arraystretch}{1.3}\setlength{\tabcolsep}{7pt}
\begin{tabular}{@{}cclll@{}}
\toprule
$N$ & $\mu$ & $\Psi_{N,\mu}(w)$ & $\sigma_j$ & $\tau_j$\\
\midrule
$5$ & $0$ & $3+3w$ & $44j^2+22j+3$ & $(4j-1)(4j-3)$\\
$5$ & $1$ & $25+35w$ & $44j^2+66j+25$ & $(4j+1)(4j+3)$\\
$7$ & $0$ & $2+6w$ & $26j^2+13j+2$ & $3(3j-1)(3j-2)$\\
$7$ & $1$ & $15+60w$ & $26j^2+39j+15$ & $3(3j+1)(3j+2)$\\
\bottomrule
\end{tabular}
\caption{The operators $L_\mu$ of~\eqref{eq:Lmu} and the recurrence
$(j+1)^2g_{j+1}=\sigma_jg_j+\tau_jg_{j-1}$ at levels $5$ and $7$. Here
$\Psi_{N,\mu}=\Psi_N+\mu\bigl(\tfrac T2+2Uw\bigr)$,
$\sigma_j=Tj(j+\mu+\tfrac12)+\Psi_{N,\mu}(0)$ and
$\tau_j=U(j-1)(j+2\mu)+\Psi_{N,\mu}'(0)$. For $\mu=0$ these are, in
characteristic zero, the recurrences of~\cite[Lemma~3.3]{Sakai14}.}
\label{tab:rec}
\end{table}

\begin{theorem}[Three-term recurrence]\label{thm:rat}
Let $N\in\{5,7\}$, $\mu=\mu_N$, and define $g_{-1}=0$, $g_0=1$ and
\begin{equation}\label{eq:ratrec}
(j+1)^2g_{j+1}=\sigma_j\,g_j+\tau_j\,g_{j-1}\qquad\text{in }\F_p,\quad 0\le j\le p-2,
\end{equation}
with $\sigma_j,\tau_j$ as in Table~\ref{tab:rec}. Then $g_j=0$ for
$n<j\le p-1$, $g_n\ne0$, and with $P_p(w)=\sum_{j=0}^{n}g_jw^j$,
\begin{equation}\label{eq:ratformula}
 S_{N,p}(X)=\ssp^{(N*)}(X)=X^{m_N}\,B_N(X)^{\mu_N}\,P_p(1/X).
\end{equation}
The polynomial $P_p$ is the unique polynomial of degree $<p$ with constant
term $1$ annihilated by $L_\mu$ modulo $p$. The polynomial
$Q_p(X)=X^{n}P_p(1/X)$ is monic, squarefree and prime to $XB_N(X)$, and
$\deg S_{N,p}=m_N+2\mu_N$.
\end{theorem}

\begin{example}\label{ex:11}
Let $N=7$ and $p=11$. Then $\mu_7=0$, $\delta_{11}=1$, $m_7=4$ and $n=3$. The
recurrence gives $(g_0,g_1,g_2,g_3)=(1,2,0,6)$ and $g_4=\dots=g_{10}=0$, so
\[
S_{7,11}(X)=X^4+2X^3+6X=X(X-2)(X^2+4X+8)\pmod{11},
\]
in agreement with $ss_{11}(J)=J(J-1728)$ and Definition~\ref{def:S}. Since
$\leg{-7}{11}=1$, this prime is in the range of Sakai's conjecture.
\end{example}

\paragraph{Heun form.}
For Heun's equation
\begin{equation}\label{eq:heunode}
\begin{split}
 z(z-1)(z-a)P''
 &+\bigl\{\gamma(z-1)(z-a)+\delta_Hz(z-a)
                    +\epsilon_Hz(z-1)\bigr\}P'\\
 &+(\alpha\beta z-q)P=0,
 \qquad \epsilon_H=\alpha+\beta+1-\gamma-\delta_H,
\end{split}
\end{equation}
the local series $\Hl(a,q;\alpha,\beta,\gamma,\delta_H;z)=\sum_jc_jz^j$ is
defined by $c_{-1}=0$, $c_0=1$ and
\begin{equation}\label{eq:heunrec}
\begin{split}
(j+1)(j+\gamma)\,a\,c_{j+1}
={}&\bigl[j\{(j-1+\gamma)(1+a)+a\delta_H+\epsilon_H\}+q\bigr]c_j\\
&-(j-1+\alpha)(j-1+\beta)c_{j-1},
\end{split}
\end{equation}
as in~\cite[p.~488]{Nakaya19}. (We write $\delta_H,\epsilon_H$ for the Heun
exponent parameters to avoid a clash with $\delta_p,\varepsilon_p$.) At level
$5$ let $\phi=(1+\sqrt5)/2$ and
$\mathcal R_p=\F_p[\phi]/(\phi^2-\phi-1)$, which is a field if
$\leg5p=-1$ and is isomorphic to $\F_p\times\F_p$ if $\leg5p=1$; at level $7$
let $\mathcal R_p=\F_p$. With the parameters of Table~\ref{tab:heun},
$\gamma=1$ and $\delta_H=\mu_N+\frac12$, the factor $(j+1)^2a$
in~\eqref{eq:heunrec} is a unit of $\mathcal R_p$ for $0\le j\le p-2$, so the
recurrence defines $c_0,\dots,c_{p-1}\in\mathcal R_p$ without ambiguity. We put
\[
\Hl_p(z)=\sum_{j=0}^{p-1}c_jz^j\in\mathcal R_p[z].
\]
This definition assumes nothing about termination.

\begin{table}[htbp]
\centering\small\renewcommand{\arraystretch}{1.4}
\begin{tabular}{@{}lll@{}}
\toprule
 & $N=5$ & $N=7$\\
\midrule
$\kappa$ & $4\phi^5$ & $27$\\
$a$ & $-\phi^{10}$ & $-27$\\
$q$ & $-(22\mu_5+3)\phi^5/4$ & $-(13\mu_7+2)$\\
$\alpha$ & $\varepsilon_p-m_5$ & $\delta_p-m_7$\\
$\beta$ & $\mu_5+\tfrac12+\tfrac14\leg{-1}{p}$ & $\mu_7+\tfrac12+\tfrac16\leg{-3}{p}$\\
$\gamma,\ \delta_H$ & $1,\ \mu_5+\tfrac12$ & $1,\ \mu_7+\tfrac12$\\
\bottomrule
\end{tabular}
\caption{The Heun parameters in Nakaya's notation, corrected. Only $\alpha$
differs from the parameters printed in Conjecture~1 of~\cite{Nakaya19}, where
it is $-m_N$. Modulo $p$ the pair $(\alpha,\beta)$ is, up to order,
$(\mu_N+\frac1r,\mu_N+1-\frac1r)$, with $r=4,3$ for $N=5,7$; for $\mu_N=0$ this
is the pair of~\cite[Conjecture~5.2]{Sakai14}.}
\label{tab:heun}
\end{table}

\begin{theorem}[Corrected Heun formulas]\label{thm:main}
Let $N\in\{5,7\}$ and let $p\ge5$ be prime, $p\ne N$. With the parameters of
Table~\ref{tab:heun},
\begin{equation}\label{eq:heunformula}
 S_{N,p}(X)=\ssp^{(N*)}(X)=X^{m_N}\,B_N(X)^{\mu_N}\,\Hl_p\!\Bigl(\frac{\kappa}{X}\Bigr).
\end{equation}
Explicitly,
\begin{align*}
\ssp^{(5*)}(X)&=X^{m_5}(X^2-44X-16)^{\mu_5}\\
&\quad\times\Hl_p\Bigl(-\phi^{10},-\tfrac{(22\mu_5+3)\phi^5}{4};\,
\varepsilon_p-m_5,\,\mu_5+\tfrac12+\tfrac14\leg{-1}{p},\,1,\,\mu_5+\tfrac12;\,
\tfrac{4\phi^5}{X}\Bigr),\\
\ssp^{(7*)}(X)&=X^{m_7}(X+1)^{\mu_7}(X-27)^{\mu_7}\\
&\quad\times\Hl_p\Bigl(-27,-(13\mu_7+2);\,
\delta_p-m_7,\,\mu_7+\tfrac12+\tfrac16\leg{-3}{p},\,1,\,\mu_7+\tfrac12;\,
\tfrac{27}{X}\Bigr).
\end{align*}
The polynomial $\Hl_p$ has degree exactly $n=m_N-o_p$, it
satisfies~\eqref{eq:heunode} over $\mathcal R_p$, and it is the only
polynomial of degree $<p$ with constant term $1$ that does. Its coefficients
satisfy $\kappa^jc_j=g_j\in\F_p$, with $g_j$ as in Theorem~\ref{thm:rat}; in
particular, at level $5$ the right side of~\eqref{eq:heunformula} lies in
$\F_p[X]$ and does not depend on the choice of $\sqrt5$. All of this remains
true, with the same polynomial $\Hl_p$, when $(\alpha,\beta)$ is replaced by
$(\mu_N+\frac1r,\,\mu_N+1-\frac1r)$.
\end{theorem}

\begin{corollary}[Sakai's Conjecture 5.2]\label{cor:sakai}
Let $N\in\{5,7\}$ and let $p\ge5$ be a prime with $\leg{-N}{p}=1$. Then
\begin{align*}
\ssp^{(5*)}(X)&=X^{m_5}\,\Hl_p\Bigl(-\phi^{10},-\tfrac{3\phi^5}{4};\,
\tfrac14,\,\tfrac34,\,1,\,\tfrac12;\,\tfrac{4\phi^5}{X}\Bigr),\\
\ssp^{(7*)}(X)&=X^{m_7}\,\Hl_p\Bigl(-27,-2;\,
\tfrac13,\,\tfrac23,\,1,\,\tfrac12;\,\tfrac{27}{X}\Bigr),
\end{align*}
where $m_N=\deg\ssp^{(N*)}$ and $\Hl_p$ is a polynomial of degree $m_N-o_p$.
\end{corollary}

This is~\cite[Conjecture~5.2]{Sakai14}, in which the level-$5$ parameters are
written $\frac{11+5\sqrt5}{11-5\sqrt5}=-\phi^{10}$,
$-\frac{3(11+5\sqrt5)}{8}=-\frac34\phi^5$ and $22+10\sqrt5=4\phi^5$, and in
which we read ``Heun polynomial over $\F_p$'' as the polynomial $\Hl_p$
defined above. The conjecture does not say that the Heun series terminates
modulo $p$, or in which degree; both are part of Theorem~\ref{thm:main}.

\begin{corollary}[The printed formula]\label{cor:printed}
Replace $\alpha$ in Table~\ref{tab:heun} by the printed value $-m_N$, generate
$c_0,\dots,c_{m_N}$ by~\eqref{eq:heunrec} in $\mathcal R_p$, and form
$X^{m_N}B_N(X)^{\mu_N}\sum_{j=0}^{m_N}c_j\kappa^jX^{-j}$. This expression
equals $S_{N,p}(X)$ if and only if $o_p=0$ or $(N,p)=(7,5)$. Thus the printed
formula holds for $p\equiv1\pmod4$ at level $5$, and for $p\equiv1\pmod3$ and
$p=5$ at level $7$; it fails for all other primes. When it fails, the
coefficient of $X^{m_N+2\mu_N-2}$ is already wrong, and at level $5$ that
coefficient does not lie in $\F_p$.
\end{corollary}

For $N=7$ and $p=11$ (Example~\ref{ex:11}) the printed parameter is
$\alpha=-4$, and the expression of Corollary~\ref{cor:printed} is
$X^4+2X^3+X^2+8X+6$, whereas $\alpha=\delta_{11}-m_7=-3$ gives
$S_{7,11}=X^4+2X^3+6X$, as do Sakai's parameters $(\frac13,\frac23)$. The
coordinate $X$ is fixed by $R_7$ and, because
$\gamma=1$, the normalized local solution at $z=0$ is unique; the discrepancy
is therefore not a matter of normalization.

\begin{corollary}[Nakaya's Conjecture 3]\label{cor:conj3}
Let $N\in\{5,7\}$ and write $f_N=\sum_{j\ge0}u_N(j)\,(j_N^*)^{-j}$. Then
$u_N(j)\in\Z$ and
\begin{equation}\label{eq:conj3}
 \ssp^{(N*)}(X)^2=B_N(X)^{\mu_N}
 \sum_{j=0}^{2(m_N+\mu_N)}u_N(j)\,X^{2(m_N+\mu_N)-j}
 \quad\text{in }\F_p[X].
\end{equation}
\end{corollary}

By the formulas of Cooper~\cite[eq.~(5.41) and Theorem~7.32]{Cooper}, as quoted
in~\cite[p.~490]{Nakaya19},
\[
u_5(j)=\binom{2j}{j}\sum_{k=0}^{j}\binom jk^2\binom{j+k}{k},\qquad
u_7(j)=\sum_{k=0}^{j}\binom jk^2\binom{2k}{j}\binom{j+k}{k},
\]
and with these closed forms~\eqref{eq:conj3} is Conjecture~3
of~\cite{Nakaya19}.

\subsection{Level 13}

Here $\Xi_{13,p}(t)=(t+6)^{\varepsilon}(t+5)^{\delta}$, $B_{13}(t)=t^2-52$,
$e=e_2+e_3$ and $n=e-\mu$. The operator
$\cL_{13,p}=\cA\,\dw^2+\cB\,\dw+\cC$ has
\[
\cA=w(1-52w^2)(1+6w)(1+5w)=w+11w^2-22w^3-572w^4-1560w^5
\]
in all cases, and $\cB$, $\cC$ as in Table~\ref{tab:ops13}. The
recurrence~\eqref{eq:rec} then expresses $(j+1)^2g_{j+1}$ through
$g_j,g_{j-1},g_{j-2},g_{j-3}$.

\begin{table}[htbp]
\centering\footnotesize
\renewcommand{\arraystretch}{1.3}\setlength{\tabcolsep}{4pt}
\begin{tabular}{@{}c l l@{}}
\toprule
$\varepsilon\delta\mu$ & $\cB$ & $\cC$\\
\midrule
000 & $1+\frac{52}{3}w-39w^2-\frac{4420}{3}w^3-4940w^4$ & $\frac{13}{6}-\frac{65}{36}w-\frac{11921}{36}w^2-\frac{8645}{6}w^3$\\
001 & $1+\frac{52}{3}w-143w^2-\frac{7852}{3}w^3-8060w^4$ & $\frac{13}{6}-\frac{3809}{36}w-\frac{64961}{36}w^2-\frac{38285}{6}w^3$\\
010 & $1+\frac{62}{3}w-19w^2-\frac{4940}{3}w^3-5980w^4$ & $\frac{23}{6}+\frac{475}{36}w-\frac{18161}{36}w^2-\frac{16445}{6}w^3$\\
011 & $1+\frac{62}{3}w-123w^2-\frac{8372}{3}w^3-9100w^4$ & $\frac{23}{6}-\frac{3269}{36}w-\frac{77441}{36}w^2-\frac{52325}{6}w^3$\\
100 & $1+\frac{70}{3}w-9w^2-\frac{5356}{3}w^3-6500w^4$ & $\frac{31}{6}+\frac{835}{36}w-\frac{23153}{36}w^2-\frac{21125}{6}w^3$\\
101 & $1+\frac{70}{3}w-113w^2-\frac{8788}{3}w^3-9620w^4$ & $\frac{31}{6}-\frac{2909}{36}w-\frac{87425}{36}w^2-\frac{60125}{6}w^3$\\
110 & $1+\frac{80}{3}w+11w^2-\frac{5876}{3}w^3-7540w^4$ & $\frac{41}{6}+\frac{1735}{36}w-\frac{29393}{36}w^2-\frac{32045}{6}w^3$\\
111 & $1+\frac{80}{3}w-93w^2-\frac{9308}{3}w^3-10660w^4$ & $\frac{41}{6}-\frac{2009}{36}w-\frac{99905}{36}w^2-\frac{77285}{6}w^3$\\
\bottomrule
\end{tabular}
\caption{The coefficients $\cB$, $\cC$ of the eight level-$13$ operators
$\cL_{13,p}$, indexed by $(\varepsilon_p,\delta_p,\mu_{13})$.}
\label{tab:ops13}
\end{table}

\begin{theorem}[Level 13]\label{thm:13}
Let $p\ge5$ be prime, $p\ne13$. Define $g_j=0$ for $j<0$, $g_0=1$ and, for
$0\le j\le p-2$,
\begin{equation}\label{eq:rec13}
(j+1)^2g_{j+1}=-\sum_{i=2}^{5}\cA_i\,(j-i+2)(j-i+1)\,g_{j-i+2}
-\sum_{i=1}^{4}\cB_i\,(j-i+1)\,g_{j-i+1}-\sum_{i=0}^{3}\cC_i\,g_{j-i}
\end{equation}
in $\F_p$, with the operator of Table~\ref{tab:ops13} that belongs to $p$.
\begin{enumerate}
\item[(i)] $g_j=0$ for $n<j\le p-1$, and $P_p(w)=\sum_{j\le n}g_jw^j$ is the
unique polynomial of degree $<p$ with constant term $1$ annihilated by
$\cL_{13,p}$ modulo $p$.
\item[(ii)] With $Q_p(t)=\sum_{j=0}^{n}g_j\,t^{n-j}$,
\begin{equation}\label{eq:13formula}
S_{13,p}(t)=(t+6)^{\varepsilon_p}\,(t+5)^{\delta_p}\,(t^2-52)^{\mu_{13}}\,Q_p(t)
\qquad\text{in }\F_p[t].
\end{equation}
The polynomial $Q_p$ is monic of degree $n$, squarefree, and prime to
$(t+6)(t+5)(t^2-52)$. Hence
\[
\deg S_{13,p}=e+\varepsilon_p+\delta_p+\mu_{13}
=\frac{7(p-1)}{12}+\frac{\varepsilon_p}2+\frac{2\delta_p}3+\mu_{13}.
\]
\item[(iii)] Let $m_0(J)=J^2-174668400\,J+10755802087387200$. If $p=5$, or if
no root of $m_0$ is a supersingular $j$-invariant in characteristic $p$, then
$g_n\ne0$ and $\deg P_p=n$. Otherwise $g_n=0$, $g_{n-1}\ne0$,
$\deg P_p=n-1$, and $t$ divides $Q_p$ exactly once.
\end{enumerate}
\end{theorem}

The roots of $m_0$ are the $j$-invariants
$87334200\pm15513120\sqrt{-13}$ of the two curves above $t=0$, that is, at
$h^2=-13$.

\begin{example}\label{ex:13}
For $p=5$ we have $(\varepsilon,\delta,\mu)=(0,1,1)$, $e_2=e_3=1$ and $n=1$.
The recurrence gives $g_1=-\cC_0=-\frac{23}{6}\equiv2$, so $Q_5=t+2$ and
$S_{13,5}=t(t^2-52)(t+2)=t^4+2t^3+3t^2+t$ modulo $5$.
For $p=149$ we have $(\varepsilon,\delta,\mu)=(0,1,1)$, $e_2=37$, $e_3=49$ and
$n=85$; here $g_{85}=0$ and $g_{84}\ne0$, so $\deg P_{149}=84$ while
$\deg Q_{149}=85$. The primes $p\le6000$ in the second case of~(iii) are
$149$, $191$, $233$, $241$, $263$, $647$, $1303$, $1709$, $1877$, $2381$,
$2593$, $2963$ and $3709$.
\end{example}

\section{Proof of the main theorem}\label{sec:proof}

In this section $N\in\{5,7,13\}$, and $h$, $K$, $A$, $x$, $w$, $d_2$, $d_3$,
$s$, $\ell_2$, $\ell_3$ are as in Section~\ref{subsec:setting}. We write
$D=q\,d/dq$, $f=f_N=(NE_2(N\tau)-E_2(\tau))/(N-1)$, and $E_2,E_4,E_6,\Delta,j$
have their usual meaning. Recall that $X_0(N)$ has genus $0$, that
$[\mathrm{SL}_2(\Z):\Gamma_0(N)]=N+1$, and that $h=q^{-1}+\cdots\in q^{-1}\Z[[q]]$
has a simple pole at the cusp $\infty$, a simple zero at the cusp $0$, and no
other zeros or poles. Consequently $x\in q^{-1}+\Z[[q]]$ and
$w\in q+q^2\Z[[q]]$, so that $\Z[[q]]=\Z[[w]]$.

\subsection{Modular identities}

\begin{lemma}\label{lem:basic}
\begin{enumerate}
\item[(a)] $Dh=-f\,h$, and $f\in1+q\Z[[q]]$ is a modular form of weight $2$ on
$\Gamma_0(N)$.
\item[(b)] For $\tau$ in the upper half-plane and even $k\ge4$,
\[
h(-1/N\tau)=\frac{K}{h(\tau)},\qquad f(-1/N\tau)=-N\tau^2f(\tau),\qquad
E_k(-1/N\tau)=(N\tau)^kE_k(N\tau).
\]
\end{enumerate}
\end{lemma}
\begin{proof}
(a) Since $D\log\eta(\tau)=E_2(\tau)/24$, we get
$D\log h=\frac{1}{N-1}\bigl(E_2(\tau)-NE_2(N\tau)\bigr)=-f$. The $q$-expansion
of $f$ is $1+\frac{24}{N-1}\sum_{n\ge1}\bigl(\sigma(n)-N\sigma(n/N)\bigr)q^n$,
with $\sigma(n/N)=0$ if $N\nmid n$, and $24/(N-1)\in\{6,4,2\}$. The modularity
of $f$ is classical.

(b) From $\eta(-1/\tau)=\sqrt{\tau/i}\,\eta(\tau)$ we get
$\eta(-1/N\tau)/\eta(-1/\tau)=\sqrt N\,\eta(N\tau)/\eta(\tau)$, whence the
first formula with $K=N^{12/(N-1)}$. From
$E_2(-1/\tau)=\tau^2E_2(\tau)+12\tau/(2\pi i)$ the non-modular terms cancel in
$NE_2(-1/\tau)-E_2(-1/N\tau)=N\tau^2\bigl(E_2(\tau)-NE_2(N\tau)\bigr)$, which
is the second formula. The third is the modularity of $E_k$.
\end{proof}

\begin{lemma}[Identity principle]\label{lem:principle}
Let $F$ be a meromorphic function on $X_0(N)$ with at most $b$ poles, counted
with multiplicity. If $F$ is holomorphic at the cusp $\infty$ and vanishes
there to order greater than $b$ in $q$, then $F=0$.
\end{lemma}
\begin{proof}
A non-zero meromorphic function on a compact Riemann surface has as many
zeros as poles.
\end{proof}

We apply the lemma to differences of quotients $g_1/g_2$ of modular forms of
equal weight $k$ on $\Gamma_0(N)$ and rational functions of $h$. By the
valence formula, $g_1/g_2$ has at most $k(N+1)/12$ poles on $X_0(N)$.

The numerators of $E_4$ and $E_6$ in terms of $h$ are the polynomials
\begin{equation}\label{eq:v46}
\begin{array}{l|l}
N=5 & v_4=h^2+250h+3125,\qquad v_6=h^2-500h-15625\\[2pt]
N=7 & v_4=h^2+245h+2401,\qquad v_6=h^4-490h^3-21609h^2-235298h-823543\\[2pt]
N=13 & v_4=h^4+247h^3+3380h^2+15379h+28561,\\
     & v_6=h^6-494h^5-20618h^4-237276h^3-1313806h^2-3712930h-4826809
\end{array}
\end{equation}
and we put $j_{\mathrm{num}}=d_3v_4^3$, a monic polynomial of degree $N+1$.

\begin{lemma}\label{lem:ids}
For $N\in\{5,7,13\}$:
\begin{enumerate}
\item[(a)] $j_{\mathrm{num}}(h)-1728h^N=d_2(h)\,v_6(h)^2$ in $\Z[h]$.
\item[(b)] $j=d_3v_4^3/h^N$, and hence $j-1728=d_2v_6^2/h^N$.
\item[(c)] $E_4=f^2\,\dfrac{v_4}{d_2d_3}$,\quad
$E_6=f^3\,\dfrac{v_6}{d_2d_3^2}$,\quad
$\Delta=f^6\,\dfrac{h^N}{d_2^3d_3^4}$.
\item[(d)] $\dfrac{2fD^2f-3(Df)^2}{4f^4}=w\,\Psi_N(w)$, with $\Psi_N$ as
in~\eqref{eq:Psi}.
\item[(e)] Put $j(h)=j_{\mathrm{num}}(h)/h^N$. Then
$R_N\bigl(J,h+A+K/h\bigr)=\bigl(J-j(h)\bigr)\bigl(J-j(K/h)\bigr)$ for the
polynomials $R_N$ of Section~\ref{subsec:setting}.
\item[(f)] $d_i(K/h)=K\,d_i(h)/h^2$ for $i=2,3$;\ \ $(h-K/h)^2=B_N(x)$; and
$w\,(h-K/h)=s(w)^{1/2}$, the square root with constant term $1$.
\end{enumerate}
\end{lemma}
\begin{proof}
(a), (e) and the first two parts of (f) are identities of polynomials or
rational functions in $h$ and are checked by expansion; in (e) one may use
$K^N(J-j(h_1))(J-j(h_2))=\Res_h\bigl(h^2-(x-A)h+K,\,Jh^N-j_{\mathrm{num}}(h)\bigr)$
for the two roots $h_1,h_2$ of $h^2-(x-A)h+K$. The last part of (f) follows
because $w(h-K/h)=(1-Ky^2)/(1+Ay+Ky^2)$ with $y=1/h$ is a power series in $w$
with constant term $1$ and square $w^2B_N(1/w)=s(w)$.

For (b), (c), (d) we use Lemma~\ref{lem:principle} for the difference of the
two sides. In (b) both sides have poles only at the two cusps, of total order
$N+1$. In (c) the quotients $E_4/f^2$, $E_6/f^3$, $\Delta/f^6$ have at most
$(N+1)/3$, $(N+1)/2$, $N+1$ poles and the rational functions at most $4$, $6$,
$N+1$. In (d), $2fD^2f-3(Df)^2$ is a multiple of the second Rankin--Cohen
bracket $[f,f]_2$, a modular form of weight $8$ on
$\Gamma_0(N)$~\cite[\S5.2]{Zagier}; so the left side has at most $2(N+1)/3$
poles, and $w\Psi_N(w)$ has at most $8$. In every case the total is at most
$2(N+1)\le28$, and the $q$-expansion of the difference vanishes to order
at least $200$ (Appendix~\ref{app:certificates}).
\end{proof}

By (b), the roots of $d_3$ and of $v_4$ are the points of $X_0(N)$ above
$j=0$, and those of $d_2$ and $v_6$ are the points above $j=1728$; the
exponents show that the roots of $d_3$ and $d_2$ are the unramified ones,
that is, the elliptic points. The parametrizations in (b) are classical; they
go back to Klein~\cite{Klein}, and~\cite{Maier} gives a modern account. For
$N=5,7$, expressions of Eisenstein series and cusp forms on $\Gamma_0^*(N)$
through $j_N^*$ and the Heun series, closely related to (c), are
in~\cite[Remark~3.2]{Sakai14}.

\subsection{The Hasse invariant}

Write $p-1=12k+4\delta_p+6\varepsilon_p$, so that $k=\lfloor p/12\rfloor$ and,
in the notation~\eqref{eq:exps},
\begin{equation}\label{eq:e2e3k}
e_2=3k+\delta_p+\varepsilon_p,\qquad e_3=4k+\delta_p+2\varepsilon_p,\qquad
6k+2\delta_p+3\varepsilon_p=\frac{p-1}{2}.
\end{equation}

\begin{lemma}[Deligne; Kaneko--Zagier]\label{lem:KZ}
$E_{p-1}\in\Zp[[q]]$ and $E_{p-1}\equiv1\pmod p$. There is a unique monic
polynomial $\tilde f\in\Zp[J]$ of degree $k$ with
\[
E_{p-1}=\Delta^k\,E_4^{\delta_p}\,E_6^{\varepsilon_p}\,\tilde f(j),
\qquad\text{and}\qquad
\ssp(J)\equiv J^{\delta_p}(J-1728)^{\varepsilon_p}\tilde f(J)\pmod p.
\]
In particular $\tilde f\bmod p$ is squarefree and
$\tilde f(0)\,\tilde f(1728)\not\equiv0\pmod p$.
\end{lemma}
\begin{proof}
The first assertion is the von Staudt--Clausen theorem. The factorization
and the congruence are the case of $E_{p-1}$ in~\cite[Theorem~1]{KZ}; this
case is due to Deligne (see~\cite{Serre}). The roots of $\ssp$ are simple by definition, and $j=0$
(resp.\ $j=1728$) is supersingular exactly when $\delta_p=1$ (resp.\
$\varepsilon_p=1$)~\cite[V.4]{Silverman}.
\end{proof}

\begin{proposition}[The Hasse quotient]\label{prop:Np}
Let $F(h)=h^{Nk}\tilde f\bigl(j_{\mathrm{num}}(h)/h^N\bigr)$ and
$N_p(h)=v_4(h)^{\delta_p}\,v_6(h)^{\varepsilon_p}\,F(h)\in\Zp[h]$. Then:
\begin{enumerate}
\item[(a)] $N_p$ is monic of degree $2e$, and
\begin{equation}\label{eq:Rp}
\frac{E_{p-1}}{f^{(p-1)/2}}=\frac{N_p(h)}{d_2(h)^{e_2}\,d_3(h)^{e_3}} .
\end{equation}
\item[(b)] The reduction $\overline N_p\in\F_p[h]$ is squarefree and prime to
$h\,d_2d_3$.
\item[(c)] For $h_0\in\Fpbar^\times$ with $d_2d_3(h_0)\ne0$, the invariant
$j(h_0)$ is supersingular if and only if $\overline N_p(h_0)=0$. At the roots
of $d_2$ we have $j=1728$, and at the roots of $d_3$ we have $j=0$.
\item[(d)] $h^{2e}\,\overline N_p(K/h)=\leg{-N}{p}K^e\,\overline N_p(h)$ in
$\F_p[h]$.
\end{enumerate}
\end{proposition}
\begin{proof}
(a) Substitute Lemma~\ref{lem:ids}(b),(c) into the factorization of
Lemma~\ref{lem:KZ} and use~\eqref{eq:e2e3k}. The degree of $N_p$ is
$(N+1)k+\delta_p\deg v_4+\varepsilon_p\deg v_6$. Comparing degrees in
Lemma~\ref{lem:ids}(c) gives $\deg v_4=\deg d_2+\deg d_3$,
$\deg v_6=\deg d_2+2\deg d_3$ and $N+1=3\deg d_2+4\deg d_3$, so
$\deg N_p=e_2\deg d_2+e_3\deg d_3=2e$.

(b) Write $\tilde f\equiv\prod_i(J-j_i)$ over $\Fpbar$ and
$G_J(h)=j_{\mathrm{num}}(h)-Jh^N$, so that $\overline F=\prod_iG_{j_i}$. One
computes
\begin{equation}\label{eq:disc}
\disc_hG_J=\pm N^{c}J^{a}(J-1728)^{b},\qquad
(c,a,b)=(65,4,2),\ (91,4,4),\ (169,8,6)
\end{equation}
for $N=5,7,13$. As $j_i\notin\{0,1728\}$ by Lemma~\ref{lem:KZ}, each $G_{j_i}$
is squarefree modulo $p$. A root $h_0$ of $G_{j_i}$ is non-zero, because
$G_J(0)=j_{\mathrm{num}}(0)$ is a power of $N$, and it satisfies
$j(h_0)=j_i$; so distinct factors $G_{j_i}$ have no common root, and
$\overline F$ is squarefree and prime to $h$. By Lemma~\ref{lem:ids}(b) the
roots of $v_4$ and $d_3$ have $j=0$ and the roots of $v_6$ and $d_2$ have
$j=1728$; these values differ modulo $p$ from each other and from every
$j_i$. Hence $v_4$, $v_6$ and $\overline F$ are pairwise coprime,
$\overline F$ is prime to $d_2d_3$, $v_4$ is prime to $d_2$, and $v_6$ is
prime to $d_3$. Finally the discriminants of $v_4$ and $v_6$, the resultants
$\Res(v_4,d_3)$ and $\Res(v_6,d_2)$, and the values $v_4(0)$, $v_6(0)$,
$d_i(0)=K$ are products of powers of $2$, $3$ and $N$
(Appendix~\ref{app:certificates}). This proves (b).

(c) By Lemmas~\ref{lem:KZ} and~\ref{lem:ids}(b),
\[
h^{N(k+\delta_p+\varepsilon_p)}\,\ssp(j(h))
\equiv d_3^{\delta_p}\,d_2^{\varepsilon_p}\,v_4^{2\delta_p}\,v_6^{\varepsilon_p}\,N_p(h)
\pmod p ,
\]
and $N_p$ is divisible by $v_4^{\delta_p}v_6^{\varepsilon_p}$. The last
sentence is Lemma~\ref{lem:ids}(b).

(d) Evaluate~\eqref{eq:Rp} at $-1/N\tau$ and use Lemma~\ref{lem:basic}(b) and
Lemma~\ref{lem:ids}(f). The denominator becomes
$K^{e}d_2^{e_2}d_3^{e_3}/h^{2e}$, and the automorphy factors combine to
$(N\tau)^{p-1}/(-N\tau^2)^{(p-1)/2}=(-N)^{(p-1)/2}$. Therefore
\[
\frac{h^{2e}\,N_p(K/h)}{K^e\,d_2^{e_2}d_3^{e_3}}
=(-N)^{(p-1)/2}\,\frac{E_{p-1}(N\tau)}{f(\tau)^{(p-1)/2}} .
\]
This and~\eqref{eq:Rp} are identities of Laurent series in $q$ with
coefficients in $\Zp$. Modulo $p$ we have
$E_{p-1}(q)\equiv E_{p-1}(q^N)\equiv1$ and
$(-N)^{(p-1)/2}\equiv\leg{-N}{p}$, so the two right sides differ by the factor
$\leg{-N}{p}$. Since $h=q^{-1}+\cdots$ is transcendental over $\F_p$, the
resulting identity of Laurent series in $q$ is an identity of polynomials in
$h$.
\end{proof}

\subsection{Descent through the Fricke involution}

\begin{lemma}[Laurent descent]\label{lem:laurent}
Let $\Bbbk$ be a field of characteristic different from $2$, let
$K\in\Bbbk^\times$ and $A\in\Bbbk$, and put $x=h+A+K/h$. Let $M\in\Bbbk[h]$ be
monic of degree $2e$ with
\[
h^{2e}M(K/h)=(-1)^\mu K^eM(h),\qquad\mu\in\{0,1\}.
\]
Then there is a unique monic $Q\in\Bbbk[x]$ of degree $e-\mu$ such that
\begin{equation}\label{eq:descent}
 M(h)=h^{e-\mu}(h^2-K)^\mu\,Q\bigl(h+A+K/h\bigr).
\end{equation}
If $M$ is squarefree, then $Q$ is squarefree and prime to $(x-A)^2-4K$. If
$c\in\Bbbk$ and $M$ is prime to $h^2+(A-c)h+K$, then $Q(c)\ne0$.
\end{lemma}
\begin{proof}
The Laurent polynomial $\Lambda=h^{-e}M(h)$ satisfies
$\Lambda(K/h)=(-1)^\mu\Lambda(h)$. The Laurent polynomials invariant under
$h\mapsto K/h$ are the polynomials in $h+K/h$, hence in $x$; an anti-invariant
one is a combination of the $h^i-(K/h)^i$, $i\ge1$, so it is $h-K/h$ times an
invariant one. This gives~\eqref{eq:descent}, and the highest exponent of
$\Lambda$ gives the degree and leading coefficient of $Q$. For
$c\in\overline{\Bbbk}$ we have $x(h)-c=(h^2+(A-c)h+K)/h$, whose two roots
are non-zero and coincide exactly when $(c-A)^2=4K$, namely at $h^2=K$. So a
repeated root of $Q$, or a root of $Q$ with $(c-A)^2=4K$, gives a repeated
root of $M$; and a root $c$ of $Q$ gives a common root of $M$ and
$h^2+(A-c)h+K$.
\end{proof}

\subsection{The differential equation}

\begin{lemma}[Uniformizing equation]\label{lem:unif}
Let $G\in\Q[[w]]$ be the series with $G(0)=1$ and $G^2=f$. Then
$G\in1+w\,\Z[\tfrac12][[w]]$ and $L_NG=0$, with $L_N$ as in~\eqref{eq:LG}.
\end{lemma}
\begin{proof}
Integrality holds because $f\in1+w\Z[[w]]$ and
$(1+u)^{1/2}\in\Z[\frac12][[u]]$. Put $y=1/h$ and $\vartheta=y\,d/dy$. By
Lemma~\ref{lem:basic}(a), $Dy=fy$, so $D=f\,\vartheta$ on $\Q[[q]]=\Q[[y]]$.
With $u=G=f^{1/2}$ this gives $\vartheta u=Du/u^2$ and
\[
\vartheta^2u=\frac{u\,D^2u-2(Du)^2}{u^5}
=\frac{2fD^2f-3(Df)^2}{4f^4}\,u=w\,\Psi_N(w)\,u
\]
by Lemma~\ref{lem:ids}(d). Next, $w=y/(1+Ay+Ky^2)$ gives
$\vartheta w=w\cdot w(h-K/h)=w\,s^{1/2}$ by Lemma~\ref{lem:ids}(f), so
$\vartheta=s^{1/2}w\,\dw$ on $\Q[[w]]$ and
\[
\vartheta^2G=w^2s\,G''+\bigl(ws+\tfrac12w^2s'\bigr)G'.
\]
Dividing $\vartheta^2G=w\Psi_NG$ by $w$ gives $L_NG=0$.
\end{proof}

\begin{lemma}[Gauge]\label{lem:gauge}
Let $(\varepsilon,\delta,\mu)\in\{0,1\}^3$ and define $\rho_2,\rho_3$
by~\eqref{eq:rhos} and $\cA,\cB,\cC$ by~\eqref{eq:Theta}
and~\eqref{eq:gaugeop}.
\begin{enumerate}
\item[(a)] $\cA,\cB,\cC\in\Z[\frac16][w]$, with $\cA(0)=0$ and
$\cA'(0)=\cB(0)=1$. For $N=5,7$ the operator is $L_\mu$ of~\eqref{eq:Lmu},
and for $N=13$ it is given by Table~\ref{tab:ops13}.
\item[(b)] The series
$Y=\ell_2^{\rho_2}\,\ell_3^{\rho_3}\,s^{-\mu/2}\,G\in1+w\,\Z[\tfrac16][[w]]$
satisfies $\cL_{N,p}Y=0$.
\end{enumerate}
\end{lemma}
\begin{proof}
(b) The factor $m=\ell_2^{-\rho_2}\ell_3^{-\rho_3}s^{\mu/2}$ has logarithmic
derivative $-\Theta$, so $G=mY$ gives $G'=m(Y'-\Theta Y)$ and
$G''=m\bigl(Y''-2\Theta Y'+(\Theta^2-\Theta')Y\bigr)$. Substituting in
$L_NG=0$ and multiplying by $\ell_2\ell_3/m$ gives $\cL_{N,p}Y=0$.
Integrality holds because $(1+u)^{r}\in\Z[\frac1c][[u]]$ for $r\in\frac1c\Z$.

(a) The statements on $\cA$ and on $\cB(0)$ are clear, since
$-2ws\Theta=\mu ws'-2ws\bigl(\rho_2\ell_2'/\ell_2+\rho_3\ell_3'/\ell_3\bigr)$
vanishes at $w=0$. For $N=5,7$ we have $\Theta=-\mu s'/(2s)$,
$\cB=s+(\mu+\frac12)ws'$ and
\[
\cC=\frac{\mu^2-\mu}{4}\,\frac{ws'^2}{s}+\frac\mu2\bigl(s'+ws''\bigr)-\Psi_N
=-\Psi_N-\mu\bigl(\tfrac T2+2Uw\bigr),
\]
because $\mu^2=\mu$ and $s'+ws''=-T-4Uw$; this is~\eqref{eq:Lmu}. For $N=13$,
$\cB$ is a polynomial for the same reason as before, and $\cC$ could have
poles only at the zeros of $s\,\ell_2\ell_3$, and of order at most one. The polar part at $s=0$ is
again a multiple of $\mu^2-\mu=0$. The residue at $\ell_2=0$ is a non-zero
multiple of $(4\rho_2+1)(4\rho_2+3)$ and the residue at $\ell_3=0$ of
$(3\rho_3+1)(3\rho_3+2)$; both vanish for the values~\eqref{eq:rhos}. (In
other words, $-\rho_2\in\{\frac14,\frac34\}$ and
$-\rho_3\in\{\frac13,\frac23\}$ are local exponents of $L_{13}$ at the elliptic
points of order $2$ and $3$.) Expanding gives Table~\ref{tab:ops13}; this
computation is certified in Appendix~\ref{app:certificates}.
\end{proof}

\begin{lemma}[Polynomial solutions modulo $p$]\label{lem:modp}
Let $\cL=\cA\,\dw^2+\cB\,\dw+\cC$ with $\cA,\cB,\cC\in\Zp[w]$, $\cA(0)=0$ and
$\cA'(0)=\cB(0)=1$.
\begin{enumerate}
\item[(a)] There is at most one polynomial $P\in\F_p[w]$ of degree $<p$ with
$P(0)=1$ and $\cL P=0$ in $\F_p[w]$. If it exists, its coefficients are the
elements $g_0,\dots,g_{p-1}$ determined by $g_0=1$ and by the vanishing
of~\eqref{eq:rec} for $0\le j\le p-2$.
\item[(b)] Let $Y\in\Zp[[w]]$ with $\cL Y=0$, and suppose that
$Y\equiv P(w)\,V(w^p)\pmod p$ with $P\in\Zp[w]$ and $V\in\Zp[[w]]$, $V(0)=1$.
Then $\cL P\equiv0\pmod p$.
\end{enumerate}
\end{lemma}
\begin{proof}
(a) The coefficient of $w^j$ in $\cL\bigl(\sum g_iw^i\bigr)$
is~\eqref{eq:rec}, in which $g_{j+1}$ has the coefficient $(j+1)^2$, a unit
for $j\le p-2$.
(b) Modulo $p$ the derivative of $V(w^p)$ vanishes, so
$0=\cL Y\equiv V(w^p)\,\cL P$, and $V(w^p)$ is a unit of $\F_p[[w]]$.
\end{proof}

Part (a) is a statement about polynomials of degree $<p$ only. In
characteristic $p$ the equation $\cL P=0$ has further formal solutions, for
instance $P_p(w)\,v(w^p)$ for any $v\in\F_p[[w]]$.

\subsection{Proof of Theorem~\ref{thm:uniform}}

\begin{proposition}\label{prop:frob}
There is a unique monic $Q_p\in\F_p[x]$ of degree $n=e-\mu_N$ with
\begin{equation}\label{eq:NpQ}
\overline N_p(h)=h^{n}\,(h^2-K)^{\mu_N}\,Q_p\bigl(h+A+K/h\bigr).
\end{equation}
It is squarefree and prime to $\xi_N(x)B_N(x)$. With $P_p(w)=w^nQ_p(1/w)$,
\begin{equation}\label{eq:frob}
\ell_2(w)^{e_2}\,\ell_3(w)^{e_3}\,s(w)^{-\mu_N/2}\,G(w)\equiv P_p(w)\,G(w^p)
\pmod p .
\end{equation}
\end{proposition}
\begin{proof}
By Proposition~\ref{prop:Np}(a),(d), Lemma~\ref{lem:laurent} applies to
$M=\overline N_p$ with $\mu=\mu_N$, since $\leg{-N}{p}=(-1)^{\mu_N}$. This
gives~\eqref{eq:NpQ}. By Proposition~\ref{prop:Np}(b), $Q_p$ is squarefree and
prime to $B_N$; and since each $d_i\ne1$ has the form $h^2+(A-c)h+K$ with $c$
the corresponding root of $\xi_N$, it is prime to $\xi_N$.

Divide~\eqref{eq:NpQ} by $d_2^{e_2}d_3^{e_3}$. As
$d_2^{e_2}d_3^{e_3}=h^e\,w^{-e}\ell_2^{e_2}\ell_3^{e_3}$,
$Q_p(x)=w^{-n}P_p(w)$ and $w(h-K/h)=s^{1/2}$, we get
\[
\frac{\overline N_p(h)}{d_2^{e_2}d_3^{e_3}}
=\frac{s^{\mu_N/2}\,P_p(w)}{\ell_2^{e_2}\ell_3^{e_3}}
\]
in $\F_p[[w]]=\F_p[[q]]$. By~\eqref{eq:Rp} and $E_{p-1}\equiv1$, the left
side is congruent to $f^{-(p-1)/2}=G^{1-p}$, and $G(w)^p\equiv G(w^p)$ modulo
$p$ because $G\in\Zp[[w]]$.
\end{proof}

\begin{proof}[Proof of Theorem~\ref{thm:uniform}]
Let $Q_p$ and $P_p$ be as in Proposition~\ref{prop:frob}; we show that they
are the polynomials of the theorem.

\emph{The differential equation.} By~\eqref{eq:rhos},
$\ell_2^{e_2}\ell_3^{e_3}=\ell_2^{\rho_2}\ell_3^{\rho_3}\cdot\ell_2^{p/4}\ell_3^{p/3}$
in $\Zp[[w]]$. For $c\in\Zp$ and $r\in\{3,4\}$ the series $(1+cw)^{p/r}$ and
$(1+cw^p)^{1/r}$ have constant term $1$ and congruent $r$th powers,
$(1+cw)^p\equiv1+cw^p$; since $r$th roots with constant term $1$ are unique
in $\F_p[[w]]$, the two series are congruent modulo $p$. Hence~\eqref{eq:frob}
gives, with $Y$ as in Lemma~\ref{lem:gauge},
\[
Y\equiv P_p(w)\,V(w^p)\pmod p,\qquad
V(w)=G(w)\,\ell_2(w)^{-1/4}\,\ell_3(w)^{-1/3}\in1+w\,\Zp[[w]].
\]
(At levels $5$ and $7$ simply $V=G$.) By Lemma~\ref{lem:gauge} and
Lemma~\ref{lem:modp}(b), $\cL_{N,p}P_p=0$ in $\F_p[w]$. Since
$\deg P_p\le n\le e_2+e_3<p$ and $P_p(0)=1$, Lemma~\ref{lem:modp}(a) shows
that the coefficients of $P_p$ are the $g_j$ of the theorem. This proves (i),
and (ii) and (iv) follow from Proposition~\ref{prop:frob} and
$Q_p(x)=x^nP_p(1/x)$.

\emph{The roots.} The polynomial $\xi_NB_N$ is squarefree modulo $p$: the
discriminant of $B_N$ is $16K$, its value at a root of $\xi_N$ is the
discriminant of the corresponding $d_i$, which is $-16$ for $d_2$ and $-27$
for $d_3$, and at level $13$ the two roots of $\xi_{13}$ differ by $1$. With
(ii), the right side of~\eqref{eq:mainformula} is therefore monic and
squarefree, and it remains to compare roots. Let $x_0\in\Fpbar$ and let
$h_0$ be a root of $h^2-(x_0-A)h+K$; the other root is $K/h_0$, and both are
non-zero. Since $\ssp$ and $R_N$ are monic in $J$, Lemma~\ref{lem:ids}(e)
shows that $x_0$ is a root of $\Res_J(\ssp,R_N)$ if and only if $j(h_0)$ or
$j(K/h_0)$ is supersingular. If $x_0$ is a root of $\xi_N$, then $h_0$ and
$K/h_0$ are the roots of $d_2$ (resp.\ $d_3$), both $j$-values are $1728$
(resp.\ $0$), and $x_0$ is a root of $S_{N,p}$ if and only if
$\varepsilon_p=1$ (resp.\ $\delta_p=1$), that is, if and only if it is a root
of $\Xi_{N,p}$. Otherwise, by Proposition~\ref{prop:Np}(c),(d), $x_0$ is a root
of $S_{N,p}$ if and only if $\overline N_p(h_0)=0$, which by~\eqref{eq:NpQ}
means $B_N(x_0)^{\mu_N}Q_p(x_0)=0$. This proves (iii).
\end{proof}

\section{Levels 5 and 7}\label{sec:57}

\begin{proof}[Proof of Theorem~\ref{thm:rat}]
By Lemma~\ref{lem:gauge}(a), $\cL_{N,p}=L_\mu$. Extracting the coefficient of
$w^j$ in $L_\mu\bigl(\sum g_iw^i\bigr)$ gives
\[
(j+1)^2g_{j+1}-\bigl(Tj(j+\mu+\tfrac12)+\Psi_{N,\mu}(0)\bigr)g_j
-\bigl(U(j-1)(j+2\mu)+\Psi_{N,\mu}'(0)\bigr)g_{j-1},
\]
which is~\eqref{eq:ratrec}. Since $\Xi_{N,p}=X^{o_p}$ and $o_p+n=m_N$,
formula~\eqref{eq:mainformula} becomes~\eqref{eq:ratformula}. Finally
$g_n=Q_p(0)\ne0$ because $Q_p$ is prime to $\xi_N=X$.
\end{proof}

\subsection{The Heun form}

\begin{proposition}\label{prop:heun}
Let $N\in\{5,7\}$, $\mu\in\{0,1\}$, and let $\kappa$, $a$ and
$q_\mu=-(22\mu+3)\phi^5/4$, resp.\ $-(13\mu+2)$, be as in
Table~\ref{tab:heun}.
\begin{enumerate}
\item[(a)] Under $z=\kappa w$, the Heun operator on the left
of~\eqref{eq:heunode} with parameters
$(a,q_\mu;\,\mu+\frac1r,\,\mu+1-\frac1r,\,1,\,\mu+\frac12)$ equals
$(a/\kappa)\,L_\mu$. At level $5$ this holds for both choices of $\sqrt5$.
\item[(b)] $G(w)=\Hl(a,q_0;\frac1r,1-\frac1r,1,\frac12;\kappa w)$ and
$s(w)^{-1/2}G(w)=\Hl(a,q_1;1+\frac1r,2-\frac1r,1,\frac32;\kappa w)$. In
particular
\[
f_N=\Hl\Bigl(a,q_0;\tfrac1r,1-\tfrac1r,1,\tfrac12;\frac{\kappa}{j_N^*}\Bigr)^2,
\]
which is~\cite[Proposition~3.1]{Sakai14}; the recurrence~\eqref{eq:ratrec}
for the coefficients of $G$ is~\cite[Lemma~3.3]{Sakai14}.
\end{enumerate}
\end{proposition}
\begin{proof}
(a) is a direct substitution, using $(\kappa w-1)(\kappa w-a)=a\,s(w)$; for
$N=5$ this identity reads
$(1-4\phi^5w)(1+4\phi^{-5}w)=1-44w-16w^2$ and holds for both conjugates of
$\phi$ (Appendix~\ref{app:certificates}). (b) The operator $L_\mu$ has
indicial polynomial $j^2$ at $w=0$, so it has a unique power-series solution
with constant term $1$. By Lemmas~\ref{lem:unif} and~\ref{lem:gauge} these
solutions are $G$ for $\mu=0$ and $s^{-1/2}G$ for $\mu=1$.
\end{proof}

The passage from $\mu=0$ to $\mu=1$ is the transformation of Heun series
in~\cite[eq.~(24)]{Nakaya19}, which is how the factor $B_N^{\mu_N}$ arises
there. The series $G$ and $s^{-1/2}G$ have rational coefficients in
$w=1/j_N^*$. The irrationality $\sqrt5$ enters at level $5$ only through the
rescaling $z=4\phi^5w$ that brings the roots of $s$ to Heun's normal position
$z=1,a$; this is why the descent to $\F_p$ and the independence of the choice
of $\sqrt5$ in Theorem~\ref{thm:main} are automatic.

\begin{proof}[Proof of Theorem~\ref{thm:main}]
Reduce Proposition~\ref{prop:heun}(a) into $\mathcal R_p$; the constants
involved, $2$, $3$, $\kappa$ and $a$, are units. The
recurrence~\eqref{eq:heunrec} depends on $\alpha,\beta$ only through
$\alpha+\beta$ and $\alpha\beta$, and by~\eqref{eq:mN} the parameters of
Table~\ref{tab:heun} satisfy
\[
\alpha=o_p-m_N\equiv\mu+o_p+\frac{\chi_p}{r},\qquad
\beta=\mu+\frac12+\chi_p\Bigl(\frac12-\frac1r\Bigr)\pmod p .
\]
So $(\alpha,\beta)\equiv(\mu+\frac1r,\mu+1-\frac1r)$ if $\chi_p=1$, and
$(\alpha,\beta)\equiv(\mu+1-\frac1r,\mu+\frac1r)$ if $\chi_p=-1$. In both
cases the Heun operator of Table~\ref{tab:heun} is $(a/\kappa)L_\mu$ modulo
$p$, so $\kappa^jc_j$ satisfies~\eqref{eq:ratrec} with the same initial
values, and $\kappa^jc_j=g_j$ for $0\le j\le p-1$. The same holds for the pair
$(\mu+\frac1r,\mu+1-\frac1r)$ itself. The theorem now follows from
Theorem~\ref{thm:rat}.
\end{proof}

\begin{proof}[Proof of Corollary~\ref{cor:sakai}]
This is the case $\mu_N=0$ of Theorem~\ref{thm:main}, with the pair
$(\frac1r,1-\frac1r)$; the degree of $\ssp^{(N*)}$ is $m_N+2\mu_N=m_N$.
\end{proof}

\begin{remark}[Origin of the discrepancy]\label{rem:typo}
Modulo $p$, the Heun operator must be the reduction of the single
characteristic-zero operator $L_\mu$, whose exponents at $z=\infty$ are
$\mu+\frac1r$ and $\mu+1-\frac1r$. Sakai's Conjecture~5.2 uses exactly these
exponents, for $\mu=0$, and is correct. In~\cite[Conjecture~1]{Nakaya19} the
pair is written $(-m_N,\beta)$, with $\beta$ depending on $\chi_p$, so as to
display a non-positive integer; this pair reduces to the correct one only when
$o_p=0$. The corrected parameter $o_p-m_N=-n$ is minus the degree of the Heun
polynomial, whereas $-m_N$ is minus the exponent of the outer power of $X$,
and the two differ exactly when the elliptic point is supersingular. The
corrected form follows the pattern of Tsutsumi's hypergeometric formulas for
$\Gamma_0(2)$ and $\Gamma_0(3)$, quoted in~\cite[Proposition~1]{Nakaya19},
which have the shape $X^{m+\varepsilon_p}\,{}_2F_1(-m,\ldots)$: the
supersingular elliptic point is carried by the outer power of $X$, and the
first parameter is minus the degree of the polynomial.
\end{remark}

\subsection{Proofs of the corollaries}

\begin{proof}[Proof of Corollary~\ref{cor:printed}]
If $o_p=0$ the printed and corrected parameters coincide. Let $o_p=1$, so
that $\chi_p=-1$ and the printed first parameter is $\alpha-1$, with $\alpha$
the corrected one. In~\eqref{eq:heunrec} this leaves $c_1=q/a$ unchanged,
lowers $\epsilon_H$ by $1$ and $\alpha\beta$ by $\beta$, so the printed $c_2$
exceeds the correct one by $(\beta-q/a)/(4a)$, where
$\beta\equiv\mu+\frac1r$. For $N=7$,
$\beta-q/a=\frac13-\frac2{27}=\frac7{27}$ if $\mu=0$ and
$\frac43-\frac{15}{27}=\frac79$ if $\mu=1$; both are units modulo
$p\ne3,7$. For $N=5$, $q/a=\frac14(22\mu+3)\phi^{-5}=\frac14(22\mu+3)(5\phi-8)$
has $\phi$-component $\frac{15}4$ or $\frac{125}4$, which is non-zero modulo
$p\ge7$, whereas $\beta\in\F_p$; so the printed $\kappa^2c_2$ does not lie in
$\F_p$. Thus, as soon as $m_N\ge2$, the printed expression differs from
$S_{N,p}$ in the coefficient of $X^{m_N+2\mu_N-2}$. With $o_p=1$, the case
$m_N\le1$ occurs only for $(N,p)=(7,5)$, where $m_7=1$ and
$c_1=\frac59\equiv0$, and both expressions equal $X(X+1)(X-27)$.
\end{proof}

\begin{proof}[Proof of Corollary~\ref{cor:conj3}]
The integrality of $u_N(j)$ follows from $f_N\in1+w\Z[[w]]$. Let
$\mathcal F(w)=\sum_ju_N(j)w^j=G(w)^2$. Squaring~\eqref{eq:frob}, in which
$\ell_2=\ell_3=1$, gives
\[
\mathcal F(w)\equiv s(w)^{\mu}\,P_p(w)^2\,\mathcal F(w^p)\pmod p .
\]
The polynomial $s^\mu P_p^2$ has degree at most $2\mu+2n=2e<p$, so it is the
truncation of $\mathcal F$ below degree $p$ modulo $p$, and
$u_N(j)\equiv0$ for $2e<j<p$. Since $B_N(X)=X^2s(1/X)$,
formula~\eqref{eq:ratformula} gives
\[
S_{N,p}(X)^2=B_N(X)^{\mu}\cdot X^{2(m_N+\mu)}\,\bigl(s^\mu P_p^2\bigr)(1/X),
\]
and $2e\le2(m_N+\mu)\le p-1$ because $m_N+\mu=(p-\chi_p)/r\le(p+1)/3$. This
is~\eqref{eq:conj3}.
\end{proof}

\begin{remark}[Lucas congruences and truncations]\label{rem:lucas}
Let $T_p$ be the truncation of $\mathcal F=\sum_ju_N(j)w^j$ below degree $p$.
The congruence $\mathcal F(w)\equiv T_p(w)\,\mathcal F(w^p)$ is equivalent to
the Lucas congruences $u_N(j_0+j_1p)\equiv u_N(j_0)\,u_N(j_1)\pmod p$ for
$0\le j_0<p$. For the sequences here these are known. Malik and
Straub~\cite{MalikStraub} proved them for the sporadic Ap\'ery-like sequences,
which include $u_7$ and the Ap\'ery numbers $\sum_k\binom jk^2\binom{j+k}{k}$
whose product with $\binom{2j}{j}$ is $u_5$; and Beukers, Tsai and
Ye~\cite{BTY} derived Lucas congruences from $E_{p-1}\equiv1$ for sequences
attached to modular forms on $\Gamma_0(N)$ and its Atkin--Lehner extensions,
by an argument close to the one used above. What Corollary~\ref{cor:conj3}
adds is the identification of the truncation: $T_p=s^{\mu}P_p^2$, it vanishes
in the degrees $2e<j<p$, and it gives the square of the supersingular
polynomial. A statement of the same type, for general genus-zero curves, is
\cite[Corollary~3]{Bogo}, where the non-ordinary locus is cut out by the
truncation of the $M$th power of the Picard--Fuchs solution, $M$ being the
least common multiple of the orders of the elliptic points; for
$\Gamma_0^*(5)$ one has $M=2$, and the square of the solution $G$ is $f_5$. We
have not checked whether the hypotheses of that result hold for the Fricke
quotients, or whether it gives~\eqref{eq:conj3} with multiplicities.
\end{remark}

\begin{remark}[The hypotheses on $p$]\label{rem:uniq}
The elements that must be units in the argument are $2$ and $3$ (square,
cube and fourth roots, and the constants in the operators), $N$ (the constant
$K$ and the discriminants in Proposition~\ref{prop:Np}), and $(j+1)^2$ for
$j\le p-2$. Together with $p\ge5$ in Lemma~\ref{lem:KZ} this is the hypothesis
$p\notin\{2,3,N\}$; no further prime is excluded. For instance, for $N=7$ and
$p=5$ we have $e=1$, $\mu_7=1$, $\overline N_5=h^2-49$, $P_5=1$ and
$S_{7,5}=X(X+1)(X-27)$.
\end{remark}

\begin{remark}[The four singular points]\label{rem:four}
The singular points of Heun's equation~\eqref{eq:heunode} are the cusp
($z=0$), the two Fricke fixed points ($z=1$ and $z=a$, the roots of $B_N$), and
the elliptic point $X=0$ ($z=\infty$), where $j=1728$ for $N=5$ and $j=0$ for
$N=7$. The elliptic point is at infinity in the variable $w$, so that it does
not enter the gauge factor in~\eqref{eq:Theta}; this is the only structural
difference from level $13$.
\end{remark}

\section{Level 13}\label{sec:13}

\begin{proof}[Proof of Theorem~\ref{thm:13}]
Parts (i) and (ii) are Theorem~\ref{thm:uniform} for $N=13$, with the
operators of Lemma~\ref{lem:gauge}(a); the second expression for the degree
follows from~\eqref{eq:exps}. For (iii), recall that $g_n=Q_p(0)$. If $p=5$,
then $t+5\equiv t$ and $Q_p(0)\ne0$, because $Q_p$ is prime to $t+5$. If
$p\ne5$, then $t=0$ is not a root of $(t+6)(t+5)(t^2-52)$ modulo $p$, and by
the proof of Theorem~\ref{thm:uniform}, $Q_p(0)=0$ if and only if $j(h_0)$ is
supersingular for $h_0^2=-13$. The two values $j(\pm\sqrt{-13})$ are the roots
of $m_0$, by a direct computation. If $Q_p(0)=0$, then $g_{n-1}\ne0$ because
$Q_p$ is squarefree.
\end{proof}

\begin{remark}[The five singular points]\label{rem:five}
In the coordinate $t$ the equation for $Q_p$ is as follows. Put
$\nu=\rho_2+\rho_3-\mu$, so that $n\equiv\nu\pmod p$, and
\begin{align*}
 \widehat{\cA}(t)&=t^5\cA(1/t)=(t+6)(t+5)(t^2-52),\\
 \widehat{\cB}(t)&=-2(\nu-1)\,t^4\cA(1/t)-t^3\cB(1/t),\\
 \widehat{\cC}(t)&=\nu(\nu-1)\,t^3\cA(1/t)+\nu\,t^2\cB(1/t)+t\,\cC(1/t).
\end{align*}
In each of the eight cases these are polynomials, of degrees $4$, $3$, $2$ and
with leading coefficients $1$, $1-2\nu$, $\nu^2$, and the substitution
$P_p(w)=w^nQ_p(1/w)$ turns $\cL_{13,p}P_p=0$ into
$\widehat{\cA}Q_p''+\widehat{\cB}Q_p'+\widehat{\cC}Q_p=0$ modulo $p$. For
example, for $(\varepsilon,\delta,\mu)=(0,0,0)$,
\[
\widehat{\cB}=\tfrac{13}{6}t^3+\tfrac{35}{2}t^2-\tfrac{92}{3}t-338,\qquad
\widehat{\cC}=\tfrac{49}{144}t^2+\tfrac{319}{144}t+\tfrac58 .
\]
In characteristic zero the operator
$\widehat{\cA}\partial_t^2+\widehat{\cB}\partial_t+\widehat{\cC}$ is Fuchsian
with the five singular points
\[
 t=-6,\qquad t=-5,\qquad t=\pm2\sqrt{13},\qquad t=\infty :
\]
the elliptic points of orders $2$ and $3$, the two Fricke fixed points, and
the cusp. These are the cusp and the four elliptic points, of orders $2,3,2,2$, of
$\Gamma_0^*(13)$, so the number of singular points is intrinsic; Heun's
equation has four. The three characters play the same roles as at levels $5$
and $7$. Cooper and Ye~\cite{CooperYe} gave a third-order equation for an
Eisenstein series of weight two and level $13$; the operator $L_{13}$ has
order two because it annihilates the square root $G$ of $f_{13}$, and we have
not found it in the literature in this form.
\end{remark}

\begin{remark}[The base point and the loss of a degree]\label{rem:basepoint}
The polynomial $Q_p$ always has degree $n$. What can drop is the degree of
its reciprocal $P_p$, and this depends on the base point $x=0$ of the
reciprocal coordinate $w=1/x$: by Theorem~\ref{thm:uniform}(iv),
$\deg P_p=n-1$ exactly when the curves above $x=0$ are supersingular and
$x=0$ is not an elliptic point. At levels $5$ and $7$ the base point is the
elliptic point, whose contribution has been removed, so $g_n\ne0$ always. At
level $13$, in the coordinate $t$, the base point carries the pair of
conjugate $j$-invariants $87334200\pm15513120\sqrt{-13}$. These are not real,
whereas singular moduli of degree two are real, so the two curves have no
complex multiplication, and the primes at which they are supersingular are
sparse (Example~\ref{ex:13}); we do not know whether there are infinitely
many. The phenomenon is not intrinsic to $\ssp^{(13*)}$. In Nakaya's
coordinate $X=t+2$ the base point is $t=-2$, above which
$j=-12288000$, the invariant with complex multiplication of discriminant
$-27$; the reciprocal polynomial formed in that coordinate would lose a degree
for every prime $p\equiv2\pmod3$. The recurrence of Theorem~\ref{thm:13}
should therefore be read as generating $Q_p$ from its leading coefficient
downwards, the last coefficient $g_n$ being allowed to vanish.
\end{remark}

\appendix
\section{Exact computational certificates}\label{app:certificates}

The proofs use finitely many characteristic-zero identities, each of which is
either an identity of polynomials or rational functions, or an identity of
meromorphic functions on $X_0(N)$ with a stated bound on the number of poles.
The supplementary PARI/GP scripts verify them exactly
(Table~\ref{tab:cert}).

\begin{table}[htbp]
\centering\small\renewcommand{\arraystretch}{1.2}
\begin{tabular}{@{}>{\raggedright\arraybackslash}p{.22\textwidth}>{\raggedright\arraybackslash}p{.46\textwidth}>{\raggedright\arraybackslash}p{.24\textwidth}@{}}
\toprule
Script & Identities & Used in\\
\midrule
\texttt{identities.gp} & Levels $5$ and $7$: the identities of
Lemma~\ref{lem:ids}; the discriminant~\eqref{eq:disc} and the auxiliary
discriminants and resultants; the gauge from $\mu=0$ to $\mu=1$; the Heun
change of variable over $\Q$ and $\Q(\sqrt5)$, for both conjugates &
Lemmas~\ref{lem:ids}, \ref{lem:gauge}; Propositions~\ref{prop:Np},
\ref{prop:heun}\\
\texttt{identities13.gp} & Level $13$: the identities of
Lemma~\ref{lem:ids}, including $R_{13}$; the discriminant~\eqref{eq:disc} and
the auxiliary discriminants and resultants; polynomiality of the eight
operators and Table~\ref{tab:ops13} &
Lemmas~\ref{lem:ids}, \ref{lem:gauge}; Proposition~\ref{prop:Np}\\
\bottomrule
\end{tabular}
\caption{Exact certificates.}
\label{tab:cert}
\end{table}

For the identities of modular functions, the scripts compare $q$-expansions
through $O(q^{200})$. The largest pole bound in Lemma~\ref{lem:ids} is $28$,
so Lemma~\ref{lem:principle} makes these comparisons conclusive. The other
identities are checked symbolically over $\Q$ or $\Q(\sqrt5)$.

Independently of the proofs, the formulas were compared with
Definition~\ref{def:S} for every admissible prime $p\le2000$, that is, for
$300$ primes at each of the three levels. These comparisons cover the
vanishing of $g_j$ for $n<j\le p-1$, the differential equations modulo $p$,
the identity~\eqref{eq:NpQ}, the squarefreeness and parity statements of
Proposition~\ref{prop:Np}, both conjugates of $\phi$ at level $5$,
Corollaries~\ref{cor:printed} and~\ref{cor:conj3}, and the criterion of
Theorem~\ref{thm:13}(iii). They are consistency checks and are not used in the
proofs.

\section*{Acknowledgments and use of AI tools}
The mathematical arguments, the computations and the text of this paper were
developed with substantial assistance from the AI systems ChatGPT (OpenAI) and
Claude (Anthropic), working under the author's direction. The literature
survey was assisted by Gemini Deep Research (Google); the statements about
\cite{Sakai14}, \cite{Bogo}, \cite{BL1} and \cite{Aricheta} were then checked
against those papers.

\begin{thebibliography}{99}

\bibitem{Aricheta}
V.~M.~Aricheta, Supersingular elliptic curves and moonshine,
\emph{SIGMA Symmetry Integrability Geom. Methods Appl.} \textbf{15} (2019),
Paper~007, 17~pp.
\href{https://doi.org/10.3842/SIGMA.2019.007}{doi:10.3842/SIGMA.2019.007}.

\bibitem{BTY}
F.~Beukers, W.-L.~Tsai and D.~Ye, Lucas congruences using modular forms,
\emph{Bull. Lond. Math. Soc.} \textbf{57} (2025), no.~1, 69--78.
\href{https://doi.org/10.1112/blms.13182}{doi:10.1112/blms.13182}.

\bibitem{Bogo}
G.~Bogo, \emph{Integrality of Picard--Fuchs differential equations of Kobayashi
geodesics and applications}, preprint (2025),
\href{https://arxiv.org/abs/2512.09188}{arXiv:2512.09188}.

\bibitem{BL1}
G.~Bogo and Y.~Li, \emph{Partial Hasse invariants for genus zero curves in
Hilbert modular varieties}, preprint (2026),
\href{https://arxiv.org/abs/2601.16929}{arXiv:2601.16929}.

\bibitem{BL2}
G.~Bogo and Y.~Li, \emph{Atkin polynomials for families of abelian varieties
with real multiplication}, preprint (2026),
\href{https://arxiv.org/abs/2601.16944}{arXiv:2601.16944}.

\bibitem{Cooper}
S.~Cooper, \emph{Ramanujan's theta functions}, Springer, 2017.

\bibitem{CooperYe}
S.~Cooper and D.~Ye, The Rogers--Ramanujan continued fraction and its level
$13$ analogue, \emph{J. Approx. Theory} \textbf{193} (2015), 99--127.

\bibitem{Dwork}
B.~Dwork, $p$-adic cycles, \emph{Inst. Hautes \'Etudes Sci. Publ. Math.}
\textbf{37} (1969), 27--115.

\bibitem{Igusa}
J.~Igusa, Class number of a definite quaternion with prime discriminant,
\emph{Proc. Nat. Acad. Sci. U.S.A.} \textbf{44} (1958), 312--314.

\bibitem{JordanNakaya}
D.~Jordan, {Nakaya’s conjectures on the linear factors and degree of Fricke supersingular polynomials}, Zenodo. (2026). \href{https://doi.org/10.5281/zenodo.24211799}{10.5281/zenodo.24211799}.

\bibitem{KZ}
M.~Kaneko and D.~Zagier, Supersingular $j$-invariants, hypergeometric series,
and Atkin's orthogonal polynomials, in \emph{Computational perspectives on
number theory}, AMS/IP Stud. Adv. Math. \textbf{7}, Amer. Math. Soc., 1998,
pp.~97--126.

\bibitem{Klein}
F.~Klein, Ueber die Transformation der elliptischen Functionen und die
Aufl\"osung der Gleichungen f\"unften Grades, \emph{Math. Ann.} \textbf{14}
(1879), 111--172.

\bibitem{Maier}
R.~S.~Maier, On rationally parametrized modular equations,
\emph{J. Ramanujan Math. Soc.} \textbf{24} (2009), no.~1, 1--73.
\href{https://arxiv.org/abs/math/0611041}{arXiv:math/0611041}.

\bibitem{MalikStraub}
A.~Malik and A.~Straub, Divisibility properties of sporadic Ap\'ery-like
numbers, \emph{Res. Number Theory} \textbf{2} (2016), Art.~5.

\bibitem{Morton21a}
P.~Morton, On the Hasse invariants of the Tate normal forms $E_5$ and $E_7$,
\emph{J. Number Theory} \textbf{218} (2021), 234--271.
\href{https://arxiv.org/abs/1906.12206}{arXiv:1906.12206}.

\bibitem{Morton21b}
P.~Morton, The Hasse invariant of the Tate normal form $E_5$ and the class
number of $\Q(\sqrt{-5l})$, \emph{J. Number Theory} \textbf{227} (2021),
94--143. \href{https://arxiv.org/abs/2002.04134}{arXiv:2002.04134}.

\bibitem{Morton23}
P.~Morton, Supersingular conjectures for the Fricke group,
\emph{Int. J. Number Theory} \textbf{19} (2023), no.~4, 803--828.
\href{https://doi.org/10.1142/S1793042123500409}{doi:10.1142/S1793042123500409}.

\bibitem{Morton24}
P.~Morton, The Hasse invariant of the Tate normal form $E_7$ and the
supersingular polynomial for the Fricke group $\Gamma_0^*(7)$,
\emph{Ramanujan J.} \textbf{63} (2024), 339--385.
\href{https://doi.org/10.1007/s11139-023-00764-8}{doi:10.1007/s11139-023-00764-8}.

\bibitem{Nakaya19}
T.~Nakaya, The number of linear factors of supersingular polynomials and
sporadic simple groups, \emph{J. Number Theory} \textbf{204} (2019), 471--496.
\href{https://doi.org/10.1016/j.jnt.2019.04.015}{doi:10.1016/j.jnt.2019.04.015}.

\bibitem{Sakai11}
Y.~Sakai, The Atkin orthogonal polynomials for the low-level Fricke groups
and their application, \emph{Int. J. Number Theory} \textbf{7} (2011), no.~6,
1637--1661.

\bibitem{Sakai14}
Y.~Sakai, The Atkin orthogonal polynomials for the Fricke groups of levels
5 and 7, \emph{Int. J. Number Theory} \textbf{10} (2014), no.~8, 2243--2255.
\href{https://doi.org/10.1142/S1793042114500766}{doi:10.1142/S1793042114500766}.

\bibitem{Serre}
J.-P.~Serre, Congruences et formes modulaires (d'apr\`es
H.~P.~F.~Swinnerton-Dyer), in \emph{S\'eminaire Bourbaki 1971/72}, Exp.~416,
Lecture Notes in Math. \textbf{317}, Springer, 1973, pp.~319--338.

\bibitem{Silverman}
J.~H.~Silverman, \emph{The arithmetic of elliptic curves}, 2nd ed.,
Graduate Texts in Mathematics \textbf{106}, Springer, 2009.

\bibitem{Tsutsumi}
H.~Tsutsumi, The Atkin orthogonal polynomials for congruence subgroups of
low levels, \emph{Ramanujan J.} \textbf{14} (2007), no.~2, 223--247.

\bibitem{Zagier}
D.~Zagier, Elliptic modular forms and their applications, in
\emph{The 1-2-3 of modular forms}, Springer, 2008, pp.~1--103.

\end{thebibliography}

\bigskip
\noindent{\scshape Maynooth University Social Sciences Institute (MUSSI), Maynooth University, Maynooth, County Kildare, Ireland}\\
\emph{Email address}: \texttt{David.Jordan@mu.ie}
\end{document}
